\begin{filecontents*}{paper_draft_reeb_space_sheet_tracking.bib}
@inproceedings{edelsbrunner2008reebspaces,
  author = {Edelsbrunner, Herbert and Harer, John and Patel, Amit K.},
  title = {Reeb Spaces of Piecewise Linear Mappings},
  booktitle = {Proceedings of the Annual Symposium on Computational Geometry},
  pages = {242--250},
  year = {2008},
  doi = {10.1145/1377676.1377720}
}

@article{tierny2017jacobi,
  author = {Tierny, Julien and Carr, Hamish},
  title = {Jacobi Fiber Surfaces for Bivariate Reeb Space Computation},
  journal = {IEEE Transactions on Visualization and Computer Graphics},
  volume = {23},
  number = {1},
  pages = {960--969},
  year = {2017},
  doi = {10.1109/TVCG.2016.2599017}
}

@article{edelsbrunner2008timevarying,
  author = {Edelsbrunner, Herbert and Harer, John and Mascarenhas, Ajith and Pascucci, Valerio and Snoeyink, Jack},
  title = {Time-Varying Reeb Graphs for Continuous Space-Time Data},
  journal = {Computational Geometry},
  volume = {41},
  number = {3},
  pages = {149--166},
  year = {2008},
  doi = {10.1016/j.comgeo.2007.11.001}
}

@article{saikia2017global,
  author = {Saikia, Himangshu and Weinkauf, Tino},
  title = {Global Feature Tracking and Similarity Estimation in Time-Dependent Scalar Fields},
  journal = {Computer Graphics Forum},
  volume = {36},
  number = {3},
  pages = {1--11},
  year = {2017},
  doi = {10.1111/cgf.13163}
}

@article{lukasczyk2020dynamic,
  author = {Lukasczyk, Jonas and Garth, Christoph and Weber, Gunther H. and Biedert, Tim and Maciejewski, Ross and Leitte, Heike},
  title = {Dynamic Nested Tracking Graphs},
  journal = {IEEE Transactions on Visualization and Computer Graphics},
  volume = {26},
  number = {1},
  pages = {249--258},
  year = {2020},
  doi = {10.1109/TVCG.2019.2934368}
}

@misc{sharma2025csp,
  author = {Sharma, Mohit and Masood, Talha Bin and List, Nanna Holmgaard and Hotz, Ingrid and Natarajan, Vijay},
  title = {Continuous Scatterplot and Image Moments for Time-Varying Bivariate Field Analysis of Electronic Structure Evolution},
  year = {2025},
  eprint = {2502.17118},
  archivePrefix = {arXiv},
  primaryClass = {cs.HC},
  doi = {10.48550/arXiv.2502.17118}
}

@misc{hristov2026singular,
  author = {Hristov, Petar and Hotz, Ingrid and Masood, Talha Bin},
  title = {Singular Arrange and Traverse Algorithm for Computing Reeb Spaces of Bivariate PL Maps},
  year = {2026},
  eprint = {2602.21087},
  archivePrefix = {arXiv},
  primaryClass = {cs.CG},
  doi = {10.48550/arXiv.2602.21087}
}
\end{filecontents*}

\documentclass[10pt,journal,compsoc]{IEEEtran}

\usepackage{amsmath}
\usepackage{amssymb}
\usepackage{booktabs}
\usepackage{array}
\usepackage{multirow}
\usepackage{graphicx}
\usepackage{xcolor}
\usepackage{url}
\usepackage[hidelinks]{hyperref}

\newcommand{\R}{\mathbb{R}}
\newcommand{\Domain}{\mathcal{M}}
\newcommand{\Range}{\mathcal{R}}
\newcommand{\Sheets}{\mathcal{S}}
\newcommand{\Sheet}{S}
\newcommand{\Score}{C}
\newcommand{\placeholder}[2]{%
  \begin{center}
  \fbox{%
    \begin{minipage}[c][#1][c]{0.92\linewidth}
      \centering
      #2
    \end{minipage}}
  \end{center}
}
\newcommand{\wideplaceholder}[2]{%
  \begin{center}
  \fbox{%
    \begin{minipage}[c][#1][c]{0.96\textwidth}
      \centering
      #2
    \end{minipage}}
  \end{center}
}

\begin{document}

\title{Tracking Reeb-Space Sheets in Time-Varying Bivariate Fields}

\author{Author~Names~Omitted~for~Draft%
\IEEEcompsocitemizethanks{
\IEEEcompsocthanksitem This is a working draft. Author names, affiliations,
acknowledgments, and final related work will be added later.}}

\markboth{IEEE Transactions on Visualization and Computer Graphics, Draft}%
{Tracking Reeb-Space Sheets in Time-Varying Bivariate Fields}

\IEEEtitleabstractindextext{%
\begin{abstract}

\end{abstract}

\begin{IEEEkeywords}
Bivariate field, Reeb space, sheet tracking, feature tracking, continuous
scatterplot, fiber surface, time-varying data, molecular visualization.
\end{IEEEkeywords}}

\maketitle



\section{Method}



\subsection{Analytics}
  \label{sec:analytics}

  A Reeb sheet is treated as a time-varying feature. When we compare sheets across consecutive timesteps, we are asking whether a sheet at time $t$ has a meaningful continuation at time $t+1$. This question has two different interpretations.

  In \emph{range mode}, a sheet continues if its geometry in the two-dimensional range space remains similar. This asks
  whether the same relationship between the two scalar fields persists.

  In \emph{domain mode}, a sheet continues if the same domain vertices continue to participate in the sheet. This asks
  whether the same spatial support persists.

  Both views are needed. A feature may look similar in range space while being supported by different domain vertices.
  Conversely, the same domain region may persist, but its range-space footprint may move or deform.

  \paragraph{Range scores.}
  For two sheets $S_t^i$ and $S_{t+1}^j$, the range score measures similarity between their range-space footprints. The
  default score is range IoU:
  \[
  q_R(S_t^i,S_{t+1}^j)
  =
  \frac{
  |R(S_t^i) \cap R(S_{t+1}^j)|
  }{
  |R(S_t^i) \cup R(S_{t+1}^j)|
  }.
  \]
  Here $R(S)$ is the sheet footprint in range space. A high value means the sheets occupy nearly the same part of range
  space. A low value means the feature has moved, deformed, disappeared, or changed its field relationship.

  We also compute auxiliary range descriptors:
  \[
  q_{\mathrm{area}} =
  \frac{\min(A_i,A_j)}{\max(A_i,A_j)},
  \]
  which checks whether the sheet sizes are similar,
  \[
  q_{\mathrm{bbox}}
  =
  \mathrm{IoU}(\mathrm{bbox}(S_t^i),\mathrm{bbox}(S_{t+1}^j)),
  \]
  which gives a coarse footprint-overlap score, and
  \[
  q_{\mathrm{centroid}}
  =
  \exp\left(
  -\frac{\|c_i-c_j\|}{D}
  \right),
  \]
  where $c_i,c_j$ are range-space centroids and $D$ is the global range-space diagonal.

  The combined range score is a weighted average:
  \[
  q_R =
  \frac{
  w_{\mathrm{iou}}q_{\mathrm{iou}}
  +w_{\mathrm{area}}q_{\mathrm{area}}
  +w_{\mathrm{bbox}}q_{\mathrm{bbox}}
  +w_{\mathrm{centroid}}q_{\mathrm{centroid}}
  }{
  w_{\mathrm{iou}}+w_{\mathrm{area}}+w_{\mathrm{bbox}}+w_{\mathrm{centroid}}
  }.
  \]
  In the current configuration, $w_{\mathrm{iou}}=1$ and the other weights are zero, so the default range score is exactly
  range IoU. The other components remain useful for experimentation, but the main analysis uses the most direct geometric
  overlap measure.

  \paragraph{Domain scores.}
  In domain mode, we compare the vertex sets of two sheets. Let $V(S)$ be the set of regular domain vertices assigned to
  sheet $S$. The raw domain overlap is
  \[
  O(S_t^i,S_{t+1}^j)
  =
  |V(S_t^i)\cap V(S_{t+1}^j)|.
  \]
  This score is meaningful because it asks whether the same spatial samples continue to belong to the feature.

  For thresholding and plotting, raw overlap counts are normalized by a dataset-level maximum so that the score lies in
  $[0,1]$. Thus, a high domain score means strong shared support, while a low domain score means that the participating
  vertices changed substantially.

  \paragraph{Threshold choice.}
  A threshold $\theta$ decides whether a match is strong enough to be treated as a continuation:
  \[
  q(S_t^i,S_{t+1}^j) \geq \theta.
  \]
  For range mode, we use fixed values such as
  \[
  \theta \in \{0.3,0.4,0.5,0.6,0.7\},
  \]
  with $0.5$ as the preferred default.

  For domain mode, the scale depends on the dataset and on vertex overlap counts. Therefore, the viewer derives domain
  thresholds at runtime. For each source sheet, it finds the best normalized domain continuation score. In domain mode, the raw score between two sheets is the number of shared domain vertices. However, this number is
  dataset-dependent: a large dataset can have much larger overlap counts than a small dataset. Therefore, before
  thresholding, we normalize the domain score.

  For each match between a source sheet $S_t^i$ and a target sheet $S_{t+1}^j$, let
  \[
  q_D(S_t^i,S_{t+1}^j)
  \]
  be the normalized domain-overlap score. In the current implementation, if the selected domain metric is the raw shared-
  vertex count, it is divided by the global maximum value of that metric in the dataset. Therefore $q_D$ lies in $[0,1]$.

  To choose a meaningful threshold, we first compute the best available continuation score for every source sheet:
  \[
  q_{\mathrm{best}}(S_t^i)
  =
  \max_j q_D(S_t^i,S_{t+1}^j).
  \]
  This is the strongest domain-overlap match that source sheet $S_t^i$ has in the next timestep. If a source sheet has no
  good continuation, its $q_{\mathrm{best}}$ will be small.

  We collect all such best scores over the dataset:
  \[
  \mathcal{Q}
  =
  \{q_{\mathrm{best}}(S_t^i)\}.
  \]
  Then we sort this list from smallest to largest. A percentile, or quantile, is simply a position in this sorted list. For
  example, the median is the middle value: half of the source sheets have $q_{\mathrm{best}}$ below it, and half have
  $q_{\mathrm{best}}$ above it.

  The viewer uses four domain threshold choices:
  \[
  \theta \in \{Q_{25}, Q_{50}, Q_{75}, Q_{90}\},
  \]
  where $Q_{25}$ is the 25th percentile, $Q_{50}$ is the median, $Q_{75}$ is the 75th percentile, and $Q_{90}$ is the 90th
  percentile of $\mathcal{Q}$. The default domain threshold is
  \[
  \theta = Q_{50}.
  \]
  This means that the default threshold is adapted to the observed domain-overlap distribution. It is not an arbitrary
  fixed number.


  \paragraph{Range intervals.}
  A range interval is the transition between two consecutive timesteps, $(t,t+1)$, analyzed using range scores. It is
  considered interesting when many sheets do not have good range-space continuations.

  For each source sheet, we find its best target sheet. For each target sheet, we also find its best source sheet. We then
  count:
  \[
  N_{\mathrm{weak}}^{src}
  =
  \#\{S_t^i : \max_j q_R(S_t^i,S_{t+1}^j)<\theta\},
  \]
  \[
  N_{\mathrm{weak}}^{tgt}
  =
  \#\{S_{t+1}^j : \max_i q_R(S_t^i,S_{t+1}^j)<\theta\}.
  \]
  The source weak count captures sheets that fail to continue. The target weak count captures sheets that appear without a
  strong predecessor.

  We also count possible splits and merges. A split occurs when one source sheet has two or more target sheets above
  threshold. A merge occurs when one target sheet has two or more source sheets above threshold. These indicate ambiguous
  or branching behavior.

  The interval event score is
  \[
  E =
  N_{\mathrm{weak}}^{src}
  +
  N_{\mathrm{weak}}^{tgt}
  +
  \lambda N_{\mathrm{split}}
  +
  \lambda N_{\mathrm{merge}}
  +
  (1-\overline{q}_{best})N_{src}.
  \]
  Here $\lambda=0.5$ in the current configuration, $\overline{q}_{best}$ is the mean best source continuation score, and
  $N_{src}$ is the number of source sheets.

  The final term is important. If every sheet barely passes the threshold, the weak counts may be zero, but the interval is
  still not stable. The continuation-gap term captures this gradual degradation.

  An interval score should capture not only missing continuations, but also structural ambiguity. A sheet may not
  disappear, but it may split into multiple plausible continuations. Similarly, multiple source sheets may merge into one
  target sheet. These are important temporal events, so they should contribute to the interval score.

  For an interval $(t,t+1)$, we count:
  \[
  N_{\mathrm{weak}}^{src}
  \]
  source sheets whose best match is below $\theta$, and
  \[
  N_{\mathrm{weak}}^{tgt}
  \]
  target sheets whose best predecessor is below $\theta$.

  We also count:
  \[
  N_{\mathrm{split}},
  \]
  the number of source sheets that have two or more target matches above $\theta$, and
  \[
  N_{\mathrm{merge}},
  \]
  the number of target sheets that have two or more source matches above $\theta$.

  The event score is
  \[
  E =
  N_{\mathrm{weak}}^{src}
  +
  N_{\mathrm{weak}}^{tgt}
  +
  \lambda N_{\mathrm{split}}
  +
  \lambda N_{\mathrm{merge}}
  +
  (1-\overline{q}_{best})N_{src}.
  \]

  Here $\lambda$ is the split/merge weight. In the current implementation,
  \[
  \lambda = 0.5.
  \]
  We use a smaller weight for splits and merges because they indicate ambiguity, but not necessarily disappearance. A weak
  or missing continuation is a stronger signal of temporal change, so it contributes weight $1$. A split or merge
  contributes weight $0.5$ by default.

  The last term,
  \[
  (1-\overline{q}_{best})N_{src},
  \]
  captures gradual weakening. Even if every sheet has a match above $\theta$, the interval may still be unstable if those
  matches are weak on average.


  The interval score also includes a term for mean continuation loss. This is needed because an interval may have no failed
  sheets, but all continuations may still be weak. For example, if \(\theta=0.5\), the best scores
  \[
  0.95,\ 0.92,\ 0.90,\ 0.88
  \]
  indicate strong continuation, while
  \[
  0.53,\ 0.52,\ 0.51,\ 0.50
  \]
  barely pass the threshold. In the second case, the weak count is still zero, but the interval is less stable.

  For an interval \((t,t+1)\), we compute the average best continuation score over all source sheets:
  \[
  \overline{q}_{\mathrm{best}}
  =
  \frac{1}{N_{\mathrm{src}}}
  \sum_i
  \max_j q(S_t^i,S_{t+1}^j).
  \]

  The gradual-weakening term is then
  \[
  (1-\overline{q}_{\mathrm{best}})N_{\mathrm{src}}.
  \]

  If the average best continuation score is high, this term is small. If the average best continuation score is low, this
  term becomes large. Therefore, it captures broad weakening of the interval even when individual sheets have not yet
  fallen below the threshold.





  \paragraph{Domain intervals.}
  Domain intervals use the same event-score idea, but replace $q_R$ with the normalized domain overlap score $q_D$. Thus, a
  high domain interval score means that the domain support is changing strongly between $t$ and $t+1$.

  Range intervals ask: ``Did the range-space features change?''

  Domain intervals ask: ``Did the spatial support of the features change?''

  These are different questions, so their high-scoring intervals need not be the same.

  \paragraph{Continuing features.}
  A continuing feature is a chain of sheet matches across time:
  \[
  S_{t_0}^{i_0}
  \rightarrow
  S_{t_0+1}^{i_1}
  \rightarrow
  \cdots
  \rightarrow
  S_{t_1}^{i_k}.
  \]
  The chain is built by keeping, for each source sheet, its best target sheet if the score is at least $\theta$. The
  feature stops when no sufficiently strong continuation exists.

  For each feature, we record its length, mean continuation score, and minimum continuation score. The length shows how
  long the feature persists. The mean score shows average stability. The minimum score shows the weakest transition along
  the chain.

  In range mode, a continuing feature means persistent range-space behavior. In domain mode, it means persistent domain
  support.

  \paragraph{Sensitivity.}
  Sensitivity analysis repeats the interval and continuing-feature computations for multiple $\theta$ values. This helps
  determine whether a result is robust.

  If the same important interval appears across many thresholds, it is likely a stable event. If it appears only for one
  threshold, it is more threshold-dependent. Similarly, if long features persist across thresholds, their continuation is
  more reliable.

  \paragraph{Domain/range complementarity.}
  Complementarity identifies places where range and domain disagree. For each source sheet, we compute its best range
  target and best domain target. If both modes choose the same target, they agree. If they choose different targets, the
  sheet has complementary behavior.

  This can happen in two important cases. First, a sheet may preserve its range-space shape but move to different domain
  vertices. Second, the same domain support may persist while its range-space behavior changes.

  For a disagreement, we measure how costly it is to force one mode to accept the other mode's target. If the range-best
  target is very different from the domain-best target, the disagreement score is high. If both alternatives are nearly
  equivalent, the score is low.

  Thus, complementarity does not simply mark all disagreements. It prioritizes disagreements where range and domain are
  confidently telling different stories.

  Domain/range complementarity identifies cases where the range and domain views suggest different continuations.

  For each source sheet $S_t^i$, we compute two best targets:
  \[
  T_R(S_t^i)
  =
  \arg\max_j q_R(S_t^i,S_{t+1}^j),
  \]
  the best range-space target, and
  \[
  T_D(S_t^i)
  =
  \arg\max_j q_D(S_t^i,S_{t+1}^j),
  \]
  the best domain-overlap target.

  If
  \[
  T_R(S_t^i)=T_D(S_t^i),
  \]
  then range and domain agree. They both say that the same target sheet is the best continuation.

  If
  \[
  T_R(S_t^i)\neq T_D(S_t^i),
  \]
  then the two modes disagree. This means that the sheet that looks most similar in range space is not the sheet that
  shares the most domain support.

  The viewer assigns a complementarity score to such disagreements. Let
  \[
  q_R^*
  =
  q_R(S_t^i,T_R(S_t^i))
  \]
  be the best range score, and let
  \[
  q_D^*
  =
  q_D(S_t^i,T_D(S_t^i))
  \]
  be the best domain score.


 For a fixed source sheet \(S_t^i\), we compare it with all candidate target sheets at the next timestep:
  \[
  S_{t+1}^1, S_{t+1}^2, \ldots, S_{t+1}^m .
  \]

  The best range-space continuation score for this source sheet is
  \[
  q_R^\ast(S_t^i)
  =
  \max_j q_R(S_t^i,S_{t+1}^j).
  \]

  The best domain-overlap continuation score for the same source sheet is
  \[
  q_D^\ast(S_t^i)
  =
  \max_j q_D(S_t^i,S_{t+1}^j).
  \]

  Thus, \(q_R^\ast\) and \(q_D^\ast\) are not global scores over the whole dataset. They are the best available range and
  domain continuation scores for one source sheet.



  We also evaluate the cross alternatives:
  \[
  q_R^D
  =
  q_R(S_t^i,T_D(S_t^i)),
  \]
  the range score of the domain-chosen target, and
  \[
  q_D^R
  =
  q_D(S_t^i,T_R(S_t^i)),
  \]
  the domain score of the range-chosen target.

  Then we compute the losses:
  \[
  L_R = \max(0, q_R^* - q_R^D),
  \]
  \[
  L_D = \max(0, q_D^* - q_D^R).
  \]

  These losses ask: how much worse is the target chosen by the other mode?

  The confidence is
  \[
  C = \min(q_R^*, q_D^*).
  \]
  This is high only when both modes have strong preferred matches. If one mode is already uncertain, the disagreement is
  less reliable.

  The complementarity score is
  \[
  C_{\mathrm{comp}}
  =
  \frac{1}{2}(L_R+L_D)C.
  \]

  A high complementarity score means that range and domain confidently choose different target sheets, and each mode
  strongly prefers its own choice over the other mode's choice. These cases are interesting because they reveal behavior
  that cannot be understood from only one view.





\section*{Acknowledgments}

Acknowledgments will be added in the final version.

\bibliographystyle{IEEEtran}
\bibliography{paper_draft_reeb_space_sheet_tracking}

\end{document}
